% Source: https://artofproblemsolving.com/community/c3h2224086_sum_of_combinations
% Stable identity: aops-topic-2224086-q01
\textbf{The sum is a single binomial coefficient.} Since $\dbinom{15}{k}=\dbinom{15}{15-k}$, the sum equals
\[\sum_{k=0}^{15}\dbinom{15}{k}\dbinom{15}{15-k}.\]
This counts the ways to choose $15$ people from $15$ boys and $15$ girls: choose $k$ boys and $15-k$ girls and add over $k$. Hence the sum is $\dbinom{30}{15}$ (Vandermonde's identity).

\textbf{Prime factorisation.} The exponent of a prime $p$ in $\dbinom{30}{15}=\dfrac{30!}{15!\,15!}$ is (exponent in $30!$) minus twice (exponent in $15!$):
\begin{itemize}
\item $p=2$: $(15+7+3+1)-2(7+3+1)=26-22=4$;
\item $p=3$: $(10+3+1)-2(5+1)=14-12=2$;
\item $p=5$: $(6+1)-2\cdot3=1$;
\item $p=7,11,13$: $4-2\cdot2=0$, $2-2\cdot1=0$, $2-2\cdot1=0$;
\item $p=17,19,23,29$: $1-0=1$ each.
\end{itemize}
So
\[\dbinom{30}{15}=2^4\cdot3^2\cdot5\cdot17\cdot19\cdot23\cdot29=720\cdot(17\cdot19\cdot23\cdot29).\]

\textbf{The last two digits.} Write $M=17\cdot19\cdot23\cdot29$. Modulo $5$, $M\equiv2\cdot4\cdot3\cdot4=96\equiv1$, so $M=5j+1$ for an integer $j$. Then
\[720M=720\cdot5j+720=3600j+720\equiv720\equiv20\pmod{100}.\]
The remainder is $\boxed{20}$. (Indeed $\dbinom{30}{15}=155117520$.)
